% Academic, restrained visual template for the Math I notes.
% Content files should keep mathematical structure only; visual styling lives here.

\usepackage{geometry}
\usepackage{fontspec}
\usepackage{unicode-math}
\usepackage{microtype}
\usepackage{enumitem}
\usepackage{xcolor}
\usepackage{booktabs}
\usepackage{titlesec}
\usepackage{fancyhdr}
\usepackage{hyperref}

\geometry{
  a4paper,
  top=2.5cm,
  bottom=2.7cm,
  inner=2.55cm,
  outer=2.25cm,
  headheight=24pt,
  headsep=0.8cm,
  footskip=1.0cm
}

\defaultfontfeatures{Ligatures=TeX}
\IfFontExistsTF{TeX Gyre Termes}
  {\setmainfont{TeX Gyre Termes}}
  {\setmainfont{Times New Roman}}
\IfFontExistsTF{TeX Gyre Heros}
  {\setsansfont{TeX Gyre Heros}[Scale=MatchLowercase]}
  {\setsansfont{Arial}[Scale=MatchLowercase]}
\IfFontExistsTF{TeX Gyre Cursor}
  {\setmonofont{TeX Gyre Cursor}[Scale=MatchLowercase]}
  {\setmonofont{Consolas}[Scale=MatchLowercase]}
\IfFontExistsTF{TeX Gyre Termes Math}
  {\setmathfont{TeX Gyre Termes Math}}
  {\setmathfont{Latin Modern Math}}

\IfFontExistsTF{SimSun}
  {\setCJKmainfont{SimSun}[AutoFakeBold=2.5, AutoFakeSlant=0.18]}
  {\setCJKmainfont{Songti SC}[AutoFakeBold=2.5, AutoFakeSlant=0.18]}
\IfFontExistsTF{Microsoft YaHei}
  {\setCJKsansfont{Microsoft YaHei}[AutoFakeBold=2.0]}
  {\setCJKsansfont{SimHei}[AutoFakeBold=2.0]}

\definecolor{aomInk}{HTML}{252522}
\definecolor{aomMuted}{HTML}{6F6A60}
\definecolor{aomRule}{HTML}{D8D1C3}
\definecolor{aomForest}{HTML}{1F3B2D}
\definecolor{aomGold}{HTML}{A9874A}
\definecolor{aomBurgundy}{HTML}{7A2E35}

\color{aomInk}
\linespread{1.18}
\setlength{\parindent}{2em}
\setlength{\parskip}{0.35em}
\setlength{\abovedisplayskip}{0.95em}
\setlength{\belowdisplayskip}{0.95em}
\setlength{\abovedisplayshortskip}{0.7em}
\setlength{\belowdisplayshortskip}{0.7em}
\numberwithin{equation}{section}
\setcounter{tocdepth}{2}
\setlist{leftmargin=2.2em, itemsep=0.3em, topsep=0.55em, parsep=0.08em}
\setlist[enumerate]{label=\arabic*., itemsep=0.3em}
\renewcommand{\arraystretch}{1.18}
\emergencystretch=2em

\hypersetup{
  colorlinks=true,
  linkcolor=aomForest,
  urlcolor=aomForest,
  citecolor=aomGold,
  filecolor=aomForest,
  pdfborder={0 0 0}
}

\titleformat{\section}
  {\sffamily\Large\bfseries\color{aomForest}}
  {\thesection}
  {0.8em}
  {}
  [\vspace{0.25ex}{\color{aomGold}\titlerule[0.45pt]}]
\titleformat{\subsection}
  {\sffamily\large\bfseries\color{aomInk}}
  {\thesubsection}
  {0.75em}
  {}
\titlespacing*{\section}{0pt}{2.6ex plus 0.8ex minus 0.2ex}{1.8ex}
\titlespacing*{\subsection}{0pt}{2.25ex plus 0.6ex minus 0.2ex}{1.1ex}

\pagestyle{fancy}
\fancyhf{}
\fancyhead[L]{\sffamily\footnotesize\color{aomMuted}考研数学一}
\fancyhead[R]{\sffamily\footnotesize\color{aomMuted}\nouppercase{\leftmark}}
\fancyfoot[C]{\sffamily\footnotesize\color{aomMuted}\thepage}
\renewcommand{\sectionmark}[1]{\markboth{\thesection\quad #1}{}}
\renewcommand{\headrulewidth}{0.35pt}
\renewcommand{\headrule}{\hbox to\headwidth{\color{aomRule}\leaders\hrule height \headrulewidth\hfill}}
\renewcommand{\footrulewidth}{0pt}

\makeatletter
\renewcommand{\maketitle}{%
  \begin{titlepage}
    \thispagestyle{empty}
    \centering
    \vspace*{0.18\textheight}
    {\sffamily\small\textls[80]{\textcolor{aomGold}{MATH I FULL-SCORE NOTES}}\par}
    \vspace{1.8em}
    {\Huge\bfseries\color{aomInk}\@title\par}
    \vspace{1.4em}
    {\color{aomGold}\rule{0.42\textwidth}{0.6pt}\par}
    \vspace{1.8em}
    {\Large\color{aomMuted}\@author\par}
    \vspace{0.8em}
    {\large\color{aomMuted}\@date\par}
    \vfill
    {\sffamily\small\color{aomMuted}高等数学 \quad 线性代数 \quad 概率论与数理统计\par}
  \end{titlepage}
}
\makeatother

\newtheoremstyle{aomplain}
  {0.8em}{0.8em}{\itshape}{}{\sffamily\bfseries\color{aomForest}}{.}{0.6em}{}
\newtheoremstyle{aomdefinition}
  {0.8em}{0.8em}{\normalfont}{}{\sffamily\bfseries\color{aomForest}}{.}{0.6em}{}
\theoremstyle{aomplain}
\newtheorem{theorem}{定理}[section]
\newtheorem{lemma}{引理}[section]
\newtheorem{corollary}{推论}[section]
\theoremstyle{aomdefinition}
\newtheorem{definition}{定义}[section]

\newcommand{\aomBlockTitle}[2]{%
  \par\addvspace{1.15\baselineskip}%
  \noindent{\sffamily\bfseries\small\color{#1}#2}\par\nobreak
  \vspace{0.4em}\nobreak
}

\newcommand{\aomBlockEnd}{%
  \par\addvspace{0.95\baselineskip}%
}

\newsavebox{\aomMetaBox}
\NewDocumentEnvironment{problemMeta}{}{%
  \par\addvspace{0.3\baselineskip}%
  \begingroup
  \footnotesize\color{aomMuted}%
  \setlength{\fboxsep}{0.65em}%
  \setlength{\fboxrule}{0.25pt}%
  \begin{lrbox}{\aomMetaBox}%
  \begin{minipage}{\dimexpr\linewidth-2\fboxsep-2\fboxrule\relax}%
  \setlength{\parindent}{0pt}%
  \setlength{\parskip}{0.18em}%
  \raggedright
}{%
  \end{minipage}%
  \end{lrbox}%
  \noindent\fcolorbox{aomRule}{white}{\usebox{\aomMetaBox}}%
  \par\endgroup\addvspace{0.6\baselineskip}%
  \ignorespacesafterend
}

\NewDocumentEnvironment{problemBox}{}{%
  \aomBlockTitle{aomGold}{题目}%
  \ignorespaces
}{%
  \aomBlockEnd\ignorespacesafterend
}

\NewDocumentEnvironment{solutionBox}{}{%
  \aomBlockTitle{aomForest}{题解}%
  \ignorespaces
}{%
  \aomBlockEnd\ignorespacesafterend
}

\NewDocumentEnvironment{knowledgeBox}{}{%
  \aomBlockTitle{aomForest}{知识点}%
  \ignorespaces
}{%
  \aomBlockEnd\ignorespacesafterend
}

\NewDocumentEnvironment{mistakeBox}{}{%
  \aomBlockTitle{aomBurgundy}{易错点}%
  \ignorespaces
}{%
  \aomBlockEnd\ignorespacesafterend
}
